\section{Short-window residual and Green means}
\label{sec:b10-short-windows}

We use the exact residual \eqref{eq:exact-relative-residual} at \(c=1\),
with physical time \(1/L_X\) held fixed throughout \(W_X\). The prime
cutoff \(P\) is constant on the outer interval \([U,U+\ell]\).
The following theorem combines the growing-cutoff diagonal, the exact
convolution weight and the short-window arm bounds
\cite{GNPRSWeightedEulerLedger,GNPRSLocalizedArm,GNPRSGrowingExponentPair,GNPRSSignedMellin}.
Write
\[
 \mathscr R(U,\ell;P)=\frac1\ell\int_U^{U+\ell}R_P(X,1)\dd X,
 \qquad
 \mu_P=\sum_{\substack{k>1\\(k,D_P)=1}}k^{-3/2}\le2P^{-1/2},
 \qquad D_P=\prod_{p\le P}p .
\]

\begin{theorem}[Short-window mean of the growing-cutoff residual]
\label{thm:b10-averaged}
Put \(\beta=17/81\). For each fixed \(0<\rho\le47/162\) and all
sufficiently large real \(U\), uniformly for integers
\(2\le P\le\lfloor L_U^2\rfloor\) and
\(U^{\beta+\rho}\le\ell\le U\),
\begin{equation}\label{eq:b10-main}
 \left|\mathscr R(U,\ell;P)-\mu_P\right|
 \le C\left(\frac{\sqrt{\mu_P}}{L_U}+L_U^{-2}\right)
     +O_\rho\!\left(U^{-\min(7/6480,\,9\rho/28)}\right).
\end{equation}
Here \(C\) depends only on the fixed inputs. The large-height onset
is finite and unevaluated.
\end{theorem}

\subsection{Exact Mellin projection of the upper part of the arm}

For the short-window proof put \(L=L_U\), \(\mathcal L=1+\log U\),
and choose
\[
 B=N_U=\left\lfloor\sqrt{U/(4\pi)}\right\rfloor,\qquad
 Z=\lfloor U^{1/3}\rfloor,\qquad Y=\lfloor U^{1/10}\rfloor,\qquad
 M=U^{1/10}.
\]
Take a fixed nonnegative smooth \(\psi_0\) of integral one supported
on \([-1/2,1/2]\) and put
\(\chi=\mathbf1_{[-3/2,5/2]}*\psi_0\).
Then \(\chi=1\) on \([-1,2]\), has support in \([-2,3]\) and
has integral four. For \(0\le a<b\le B\) define
\begin{align*}
 w_x=w_x(U)&=x^{-3/4}\exp((\log x)^2/(4L)),\qquad
 \tau(u)=u/2+\pi/(8L),\\
 D_{a,b}(u)&=\sum_{a<n\le b}w_n\,\ee^{\ii\tau(u)\log n},\\
 \mathcal Q_{a,b}&=\ell^{-1}\int_\R
       \chi((u-U)/\ell)|D_{a,b}(u)|^2\dd u,\\
 \mathcal S_{a,b}&=\sum_{a<m<n\le b}w_mw_n
       \ee^{\ii\tau(U)\log(n/m)}
       \widehat\chi\bigl((\ell/2)\log(n/m)\bigr),
 \qquad
 \widehat\chi(\xi)=\int_\R\chi(v)\ee^{\ii v\xi}\dd v.
\end{align*}
Expanding the square and pairing the two orderings gives
\begin{equation}\label{eq:b10-signed}
 \mathcal Q_{a,b}=4\sum_{a<n\le b}w_n^2+2\Re\mathcal S_{a,b}.
\end{equation}

We first show, for \(1\le\ell\le\sqrt U\), that
\begin{equation}\label{eq:b10-high}
 \mathcal Q_{Z,B}+|\Re\mathcal S_{Z,B}|
 \ll\mathcal L^3\left(U^{-1/9}+\frac{U^{37/180}}{\ell}\right).
\end{equation}
Use the cumulative zeta mean-square error of Bourgain and Watt
\cite[Theorem 4]{BourgainWattPerturbedCones2015v1}, published in
\cite{BourgainWattPerturbedCones2018}. For \(T\ge1\) they give
\[
 \int_0^T|\zeta(1/2+\ii t)|^2\dd t
 =T\bigl(\log(T/(2\pi))+2\gamma-1\bigr)
      +O_\epsilon(T^{1515/4816+\epsilon}).
\]
Fix \(\theta=19/60\) and
\(\epsilon=151/144480<\theta-1515/4816\).
For a subinterval of \([1,R]\) subtract the two cumulative
identities, integrating the derivative
\(\log(t/(2\pi))+2\gamma\) of the main term.
The two endpoint errors cost \(O(R^\theta)\). Reflect negative heights
and use boundedness on \([-1,1]\). For every interval \(I\subset[-R,R]\)
of length at most \(c_*\ell\), \(R\ge2\), \(\ell\ge1\), this gives
\begin{equation}\label{eq:b10-local-zeta}
 \int_I|\zeta(1/2+\ii t)|^2\dd t
 \ll_{c_*} \ell\log(2+R)+R^\theta.
\end{equation}
In particular the cumulative error remains in the translated bound.

Let \(q\) be a fixed smooth increasing cutoff with \(q=0\) on
\((-\infty,0]\) and \(q=1\) on \([1/2,\infty)\), and set
\[
 a(x)=\sqrt x\,w_x,\qquad
 b(x)=a(x)q(x-Z-1/4)q(B+3/4-x),\qquad
 \widetilde b(z)=\int_0^\infty b(x)x^z\frac{\dd x}{x}.
\]
The transitions lie in \((Z+1/4,Z+3/4)\) and
\((B+1/4,B+3/4)\). At every integer \(n\),
\(b(n)=\sqrt n\,w_n\mathbf1_{Z<n\le B}\).
Put \(A_*=U^{-1/18}\) and \(g(y)=b(\ee^y)\). On the support,
\(\alpha=\log x/\log U\) ranges from \(1/3+O(U^{-1/3}/\log U)\)
to \(1/2+O(1/\log U)\).
Since \(\log a(U^\alpha)=-\alpha(1-\alpha)\log U/4+O(1)\),
we have \(\sup a\ll A_*\). The logarithmic derivative of \(a(\ee^y)\)
with respect to \(y\) is \(-1/4+y/(2L)\) and changes sign at most once.
The interior total variation is \(O(A_*)\) and the interior
second derivative has integral \(O(A_*\mathcal L)\).
A transition at \(x\asymp N\), \(N=Z,B\), has \(y\)-width
\(O(1/N)\) and first and second \(y\)-derivatives bounded by
\(O(N)\) and \(O(N^2)\). The product rule gives
\[
 \|g\|_1\ll A_*\mathcal L,\qquad
 \|g'\|_1\ll A_*,\qquad
 \|g''\|_1\ll A_*B.
\]
For \(k(v)=\widetilde b(\ii v)/(2\pi)\) integration by parts yields
\[
 |k(v)|\ll A_*\min\{\mathcal L,|v|^{-1},B|v|^{-2}\}
 \quad(v\ne0).
\]
Split the integrals at \(1/\mathcal L,B,U\), with \(B\le\sqrt U\).
The resulting bounds are
\begin{align*}
 \int_\R |k(v)|\dd v&\ll A_*\mathcal L,\\
 \int_\R |k(v)|\log(2+U+|v|)\dd v&\ll A_*\mathcal L^2,\\
 \int_\R |k(v)|(U+|v|)^\theta\dd v&\ll A_*U^\theta\mathcal L.
\end{align*}
For the last two integrals, the \(|v|>U\) pieces cost
\(O(A_*B(1+\log U)/U)\) and \(O(A_*B U^{\theta-1}/(1-\theta))\),
respectively.

With \(s=1/2-\ii\tau(u)\), Mellin inversion first on \(\Re z=1\)
permits the absolutely convergent Dirichlet series for
\(\zeta(s+z)\). Shift to \(\Re z=0\). For each fixed \(U,u\), smooth
compact support gives arbitrary vertical decay of \(\widetilde b\) on
the contour strip, and Euler summation gives polynomial growth of
\(\zeta\) away from its pole. The crossed simple pole has residue
\(\widetilde b(1-s)\). Thus the exact identity is
\begin{equation}\label{eq:b10-mellin}
 D_{Z,B}(u)=\sum_n b(n)n^{-s}
   =\widetilde b(1-s)+\int_\R k(v)\zeta(s+\ii v)\dd v .
\end{equation}
The weight \(\ee^{y/2}a(\ee^y)\) increases on the support and is bounded
by \(C A_*\sqrt B\). After multiplying by the cutoffs its variation is
still \(O(A_*\sqrt B)\), so one integration by parts pays the residue
by \(C A_*\sqrt B/|\tau(u)|\ll A_*U^{-3/4}\) on the \(\chi\)-window.

Apply Cauchy--Schwarz to the last integral in \eqref{eq:b10-mellin}
with measure \(|k(v)|\dd v\), then use Tonelli's theorem with
\(\chi((u-U)/\ell)\dd u/\ell\). For each \(v\), the substitution
\(t=v-\tau(u)\) has absolute Jacobian \(|\dd u/\dd t|=2\).
It sends the support into an interval of length \(5\ell/2\)
in \([-C(U+|v|),C(U+|v|)]\).
Equation~\eqref{eq:b10-local-zeta} and the weighted kernel bounds give
\[
 \mathcal Q_{Z,B}
 \ll A_*^2\left(U^{-3/2}+\mathcal L^3+
                      \frac{U^\theta\mathcal L^2}{\ell}\right).
\]
The diagonal satisfies \(\sum_{Z<n\le B}w_n^2\ll A_*^2\mathcal L\).
Since \(\theta-1/9=37/180\), \eqref{eq:b10-signed} proves
\eqref{eq:b10-high}. Its signed estimate concerns the real quadratic
sum over all pairs, including all lattice resonances.

\subsection{Lower modes and the weighted arm norms}
For the shorter windows \(\ell\le\sqrt U\), use the auxiliary arm
\(A(u)\) at time \(1/L_u\) from \eqref{eq:real-two-arms}, with
\(A_Y(u)\) its restriction to \(n\le Y\). Its weights satisfy
\(w_n(u)\le n^{-5/8}\) for \(n\le N_u\). Define
\begin{align*}
 r_*&=\frac{32}{\pi}\log(2U),&
 \mathcal I&=[U-r_*,U+\ell+r_*],\\
 W(u)&=\int_U^{U+\ell}K(u-X)\mathbf1_{u\in W_X}\dd X,&
 d\sigma(u)&=\frac{2W(u)}{\ell}\dd u .
\end{align*}
The B07 weight \cite{GNPRSLocalizedArm} has \(0\le W\le1/2\),
support in \(\mathcal I\), and
\begin{equation}\label{eq:b10-weight}
  1-\frac4\pi U^{-4}\le m_\sigma:=\int d\sigma\le1,\qquad
  \operatorname{TV}(W)\le1+\frac{\ell}{2U^4}.
\end{equation}
The mass bound integrates the kernel tail
\(\int_{|v|>R_X}K(v)\dd v\le (2/\pi)X^{-4}\).
The untruncated convolution has derivative
\(K(u-U)-K(u-U-\ell)\) of \(L^1\) norm at most one.
Each removed tail has total variation \(4K(R_X)\le X^{-4}/2\),
including the two jumps. Integrating their variations proves the
second bound. For \(r_*\le\ell\), \(\mathcal I\subset[U-\ell,U+2\ell]\),
where the smooth weight \(\chi((u-U)/\ell)\) equals one. Thus it
majorizes \(2W(u)\) in upper bounds for nonnegative norm squares.

On the full smooth support \([U-2\ell,U+3\ell]\), freeze the
weights and the small \(\pi/(8L_u)\) phase correction at \(U\).
For common active indices their summed drift is
\(O(\ell U^{-13/16})\), using \(O(U^{-1})\) bounds for the
logarithmic weight derivative and for the derivative of the
small phase after multiplication by \(\log n\). At the upper
cutoff there are \(O(1+\ell/\sqrt U)\) changed indices of size
\(\asymp\sqrt U\), each of weight \(O(U^{-5/16})\).
With \(Y\le N_{U-2\sqrt U}\), this gives uniformly
\[
 |(A-A_Y)-D_{Y,B}|+|A_Y-D_{0,Y}|\ll U^{-5/16}.
\]
The cutoff jumps therefore cost \(O(U^{-5/8})\) in the norm-square
comparison, which uses \(|f+g|^2\le2|f|^2+2|g|^2\).

In the Fourier expansion of \(\mathcal Q_{Y,Z}\), partition the
pairs \(Y<m<n\le Z\) into \(n\ge2m\) and \(m<n<2m\).
For the latter write \(n=m+h\) and partition \(m\) into
dyadic blocks \(N\le m<2N\). For fixed \(h\) the cutoff
conditions restrict \(m\) to an interval in this block. Put
\(t=(U/2+\pi/(8L))/(2\pi)\asymp U\). The conjugate exponential
phase, in the convention \(\exp(2\pi\ii f)\), is
\[
 f(x)=-t\log(1+h/x)=-t\int_0^h(x+s)^{-1}\dd s .
\]
For \(\varphi(x)=-th/x\), every derivative order \(j\ge1\) satisfies
\[
 \frac{|f^{(j)}(x)-\varphi^{(j)}(x)|}{|\varphi^{(j)}(x)|}
 \le \frac{(j+1)h}{2x}.
\]
This follows by averaging
\(1-(1+s/x)^{-j-1}\le(j+1)s/x\). Take the derivative
parameter \(2\) and scale \(th\) in the exponent-pair class
of Montgomery and Vaughan \cite[Definitions 16.1--16.2 and
Theorems 16.23, 16.25]{MontgomeryVaughanExponentialSums2026}.
The operations \(B,A,A,B\) from \((0,1)\) give \((2/7,4/7)\).
Fix the derivative order and tolerance for this pair, then
choose a fixed \(\delta_0>0\) small enough that \(h\le\delta_0N\)
ensures the required relative derivative bounds. The unweighted
sum over each subinterval is
\[
 \ll (th/N^2)^{2/7}N^{4/7}+N^2/(th)
 \ll (Uh)^{2/7}+N^2/(Uh).
\]
Both terms are used.

The exact weight is decreasing in its index, with
\(|w_x'(U)|\ll w_x(U)/x\). For
\[
 a_h(x)=w_x(U)w_{x+h}(U)
          \widehat\chi\bigl((\ell/2)\log(1+h/x)\bigr)
\]
the supremum plus total variation on each allowed interval is
\(\ll_D w_N(U)^2(1+\ell h/N)^{-D}\) for fixed \(D\).
Indeed the weight derivatives cost \(O(w_N(U)^2/N)\),
and one extra Fourier decay power pays the integral of the
argument derivative \(O(\ell h/N^2)\).
Abel summation and
\(\sum_{h\ge1}h^{2/7}(1+\ell h/N)^{-D}\ll (N/\ell)^{9/7}\)
for \(D\ge4\) give the leading cost per block
\[
 U^{2/7}N^{-3/14}\exp\bigl((\log N)^2/(2L)\bigr)\ell^{-9/7}.
\]
Writing \(\alpha=(\log N)/(\log U)\), the exponent of \(U\)
in its coefficient is \(2/7-3\alpha/14+\alpha^2/2+O(1/\log U)\).
The quadratic is convex. Its endpoint values on \([1/10,1/3]\)
are \(377/1400\) and \(17/63\), with difference \(1/1800\).
Floor and dyadic factors cost \(O(1)\) in the logarithm of the
coefficient. The block count is \(O(\mathcal L)\).
For the \(N^2/(Uh)\) term, the coarser \(N^{-5/4}\) amplitude
and \(\sum_{h\ge1} h^{-1}(1+\ell h/N)^{-D}\ll\log(2+N/\ell)\)
give \(O(U^{-1}Z^{3/4}\mathcal L)\) after dyadic summation.
For \(h>\delta_0 N\) and for noncomparable pairs, the log ratio
is bounded below by a fixed positive constant. The Fourier
factor is \(O(\ell^{-4})\) and the squared absolute sum of the
weights is \(O(U^{3/8})\). The diagonal is at most \(16Y^{-1/4}\).
Altogether,
\begin{equation}\label{eq:b10-low}
 \mathcal Q_{Y,Z}\ll Y^{-1/4}+\mathcal L U^{17/63}\ell^{-9/7}
                       +U^{3/8}\ell^{-4}+U^{-5/8}\mathcal L .
\end{equation}

The exact split \(D_{Y,B}=D_{Y,Z}+D_{Z,B}\) pays all cross
terms through the square triangle bound with factor two.
In the frozen head the diagonal is bounded, and the distinct
log spacings are at least \(1/Y\). Second-power Fourier decay
and \((\sum_{n\le Y}w_n(U))^2\ll Y^{3/4}\) give the
off-diagonal cost \(O(Y^{11/4}/\ell^2)\), which decays on the
short-window range in Theorem~\ref{thm:b10-averaged}.
Together with \eqref{eq:b10-high}, \eqref{eq:b10-low}, and the
freezing comparison this gives, for a sufficiently large fixed \(C\),
\begin{align}
 \mathcal E&=\mathcal L U^{17/63}\ell^{-9/7}
      +\mathcal L^3(U^{-1/9}+U^{37/180}/\ell)
      +U^{3/8}\ell^{-4}+U^{-5/8}\mathcal L,\nonumber\\
 \int |A-A_Y|^2d\sigma&\le A_1:=C(Y^{-1/4}+\mathcal E),\nonumber\\
 \int |A|^2d\sigma,\quad\int |A_Y|^2d\sigma
   &\le A_0:=C(1+\mathcal E).
 \label{eq:b10-norms}
\end{align}

\subsection{The residual estimate}
At auxiliary time \(1/L_u\) write \(Q=E_P^{-1}\), \(D_P=\prod_{p\le P}p\),
and for an arm \(a\) put
\[
 \mathcal Z(a)=aQ-1+\ee^{\ii\phi}(Q/\overline Q)\overline{aQ-1},
 \qquad
 \begin{aligned}
 \phi(u)&=u(L_u-1)/2+\pi/8+7\pi/4,\\
 \nu(u)&=u/2+\pi/(8L_u).
 \end{aligned}
\]
The weighted residual estimate uses the B06 coefficient and clock
bounds \cite{GNPRSWeightedEulerLedger} and B07's exact boundary
calculation. First compare the physical square at time \(1/L_X\),
fixed throughout \(W_X\), to \(|\mathcal Z(A)(u)|^2\).
At fixed \(u\in W_X\) the two auxiliary arms have size \(O(U^{7/32})\)
and time derivative \(O(L^2U^{7/32})\) on the segment from \(1/L_X\)
to \(1/L_u\), by the weights and carrier phase in
\eqref{eq:real-two-arms}. On this segment keep the Euler exponent
\(3/4\). With
\[
 H=\prod_{p\le P}(1-p^{-3/4})^{-1},\qquad
 T_1(P)=\sum_{p\le P}\frac{p^{-3/4}\log p}{1-p^{-3/4}},
\]
the reciprocal product has modulus at most \(H\) and time derivative
\(O(HT_1(P))\). Its ratio with its conjugate has unit modulus and
time derivative \(O(T_1(P))\). Since \(T_1(P)\ll P^{1/4}\log P\le L^2\)
eventually, the relative square has time derivative
\(O(H^2L^2 U^{7/16})\). The displacement
\(O(r_* /(UL^2))\) costs \(O(H^2r_*U^{-9/16})\).
The exact field error in \eqref{eq:real-provider-error}, after
multiplication by \(Q\) and comparison of squares, adds
\(O(H^2U^{-1/32})\). This calculation differentiates only the
arithmetic approximation at fixed \(u\). Fubini now gives
\begin{equation}\label{eq:b10-fubini}
 \left|\mathscr R-\frac12\int|\mathcal Z(A)|^2\,d\sigma\right|
 \le C_{\rm clk}H^2(U^{-1/32}+r_*U^{-9/16}).
\end{equation}
The factor \(1/2\) comes from the definition of \(d\sigma\).

Let \(Q_M,E_M\) denote the expansions of \(Q,E_P\) truncated
by integer index at \(M\). Their indices are squarefree divisors
of \(D_P\) for \(Q_M\) and integers with all prime factors at most
\(P\) for \(E_M\). Their absolute coefficient sums are at most
\(H\), as for the complete expansions. Put \(B_M=A_YQ_M-1\) and
\(m_M=\overline{Q_M}E_M\). With
\[
 T_M=M^{-1/8}\prod_{p\le P}(1-p^{-5/8})^{-1},
\]
both discarded tails have supremum at most \(T_M\). This follows
from \(n^{-3/4}\le M^{-1/8}n^{-5/8}\) for \(n>M\).
Also \(\|m_M-\overline Q/Q\|_\infty\le2HT_M\), and
the coefficient absolute sum of \(m_M\) is at most \(H^2\).

For \(L_0=L_{U-r_*}\), \(K_0=\lfloor M\rfloor Y\) set
\begin{align*}
 a_0&=C_2/(4L_0)+C_0Y^{-1/8}+C_0M^{-1/8},
 &D_0&=2\sqrt{\mu_P}\,a_0+a_0^2,\\
 s_0&=150000 K_0(1+\log K_0)/\ell,
 &b_0&=(\sqrt{\mu_P}+a_0)^2+s_0,\\
 c_0&=8000 H^4Y^{3/4}/(\ell L_0).
\end{align*}
The coefficient vector of \(B_M\) has distance at most \(a_0\)
in \(\ell^2\) from \(k^{-3/4}\mathbf1_{k>1,(k,D_P)=1}\),
by Möbius cancellation. More explicitly, the individual difference
is bounded by \(k^{-5/8}d(k)\) times
\((\log k)^2/(4L_0)+Y^{-1/8}+M^{-1/8}\), where \(d(k)\) counts
divisors. The three terms bound the Gaussian weight correction,
the head truncation and the reciprocal truncation, respectively.
Here \(C_0^2=\sum_k d(k)^2 k^{-5/4}\) and
\(C_2^2=\sum_k d(k)^2(\log k)^4k^{-5/4}\), which converge by
\(d(k)^2\le d_4(k)\) and the absolute Dirichlet series for
\(\zeta(s)^4\), \(s>1\). In particular \(C_0\le25\).
The diagonal error at each \(u\) is at most \(D_0\).

For same-arm phases \(\nu\log(k/l)\), \(k\ne l\), the partial
integrals are at most \(12/|\log(k/l)|\) by monotonicity
and \(1/4\le\nu'\le1/2\). The coefficients have supremum at most
\(g_k=k^{-5/8}d(k)\) and variation at most \(2g_k\) by summing
the decreasing weights over divisors, including any activation
jumps. For the product of two coefficients multiplied by \(W\),
the supremum plus variation is at most \(4g_kg_l\) by
\eqref{eq:b10-weight}. Using
\(|\log(k/l)|^{-1}\le K_0/|k-l|\), \(2g_kg_l\le g_k^2+g_l^2\)
and the harmonic row sum gives the weighted off-diagonal cost
\(s_0\), with the factor \(2/\ell\) from \(d\sigma\). Hence
\[
 \left|\int|B_M|^2d\sigma-\mu_P\right|
 \le D_0+s_0+\mu_P(1-m_\sigma),\qquad
 \int|B_M|^2d\sigma\le b_0.
\]
The reflected phases are \(-\phi+\nu\omega\) with
\(\omega=\log(kle/d)\), \(k,l\le K_0\), \(e,d\le M\).
Their absolute log frequencies are bounded by
\(2\log Y+3\log M\le\tfrac12\log U\).
On \(\mathcal I\), eventually this is at most \(7L_0/8\),
and \(\nu''|\omega|\ll U^{-2}/L\le1/(12U)\) while
\(\phi''=1/(2u)\ge1/(6U)\). The first derivative is thus
monotone and at most \(-L_0/16\), giving partial integrals
bounded by \(48/L_0\).
The pre-convolution absolute coefficient envelope of \(B_M\)
sums to at most \((8/3)HY^{3/8}\), since the coefficient at
1 cancels. Combining the product variation estimate for these
terms with the absolute coefficient sum \(H^2\) of \(m_M\)
bounds \(\big|\int 2\Re(\ee^{-\ii\phi}B_M^2m_M)d\sigma\big|\)
by \(c_0\).

The exact square for \(A_Y\) is
\(2|A_YQ-1|^2+2\Re(\ee^{-\ii\phi}(A_YQ-1)^2\overline Q/Q)\).
To replace the complete arm by \(A_Y\), use
\(\|\mathcal Z(A)-\mathcal Z(A_Y)\|_{L^2(d\sigma)}
 \le2H\sqrt{A_1}\) and
\(\|\mathcal Z(A)\|_{L^2(d\sigma)},\|\mathcal Z(A_Y)\|_{L^2(d\sigma)}
 \le4H\sqrt{A_0}\). For the multiplier replacement use
\(\|(A_YQ-1)-B_M\|_{L^2(d\sigma)}\le T_M\sqrt{A_0}\),
the bound \(b_0\), and \(|\overline Q/Q|=1\).
Combining the square differences, the diagonal and reflected
bounds, and \eqref{eq:b10-fubini}, yields
\begin{align}
 |\mathscr R-\mu_P|\le{}&
 D_0+s_0+\tfrac12c_0+8H^2\sqrt{A_0A_1}
    +4T_M\sqrt{A_0b_0}+2T_M^2A_0+2HT_M b_0\nonumber\\
 &+C_{\rm clk}H^2(U^{-1/32}+r_*U^{-9/16})
    +\frac4\pi\mu_P U^{-4}.
 \label{eq:b10-ledger}
\end{align}

To conclude Theorem~\ref{thm:b10-averaged}, take
\(\beta=17/81\) and \(\ell\ge U^{\beta+\rho}\).
In \eqref{eq:b10-norms} the low term has power \(-9\rho/7\),
and the translated high term has power \(-7/1620-\rho\),
before logarithms. The diagonal tail has power \(-1/40\),
the remaining large-shift cost has power \(-301/648-4\rho\),
and freezing and the phase remainder cost \(O(U^{-5/8}\mathcal L)\).
The head's auxiliary off-diagonal has power \(-469/3240-2\rho\).
In \eqref{eq:b10-ledger}, \(s_0\ll\mathcal L U^{-4/405-\rho}\)
and \(c_0\ll H^4U^{-437/3240-\rho}/L_0\).
The two coefficient tails in \(a_0\) and the power in \(T_M\)
are \(O(U^{-1/80})\). In particular \(A_0,b_0=O(1)\) eventually.
Integral comparison gives \(H\le\exp(CP^{1/4})\) and
\(T_M\le U^{-1/80}\exp(CP^{3/8})\).
The remaining polynomial errors, including the field, clock and
mass terms, fit
\[
 |\mathscr R-\mu_P|
 \le C(\sqrt{\mu_P}/L+L^{-2})
 +C\exp\bigl(C(P^{1/4}+P^{3/8})\bigr)\mathcal L^{C'}
       (U^{-7/3240}+U^{-9\rho/14}).
\]
The positive exponents here precede absorption of the product and
logarithm factors. Uniformly for \(P\le L^2\), the logarithm of
those factors is \(O(L^{3/4}+\log\mathcal L)=o(\log U)\). It
can be bounded by \(\tfrac12\min(7/3240,9\rho/14)\log U\)
after a finite onset. This proves the stated remainder. All
preceding gates hold at finite heights for fixed \(\rho>0\):
\(r_*\le\ell\), \(\mathcal I\subset[U/2,3U]\), \(L_0\ge16\),
\(2\le Y<Z<B\), \(Y\le N_{U-2\sqrt U}\), and the field,
frequency and product bounds above. The common onset is
unevaluated. For \(\sqrt U\le\ell\le U\), the inherited B08
translated residual theorem \cite[Eq.~(B08.1)]{GNPRSGrowingExponentPair}
with its fixed margin \(19/72\) starts at \(\ell=\sqrt U\)
and gives the same leading logarithmic terms with remainder
\(O(U^{-1/144})\). It uses the same block-constant \(P\) and
physical clock.

\subsection{Green consequences at the block cutoff}
\begin{corollary}\label{cor:b10-green}
On the interval range of Theorem~\ref{thm:b10-averaged}, take
\(P_U=\lfloor L_U^2\rfloor\), constant over \([U,U+\ell]\).
Then
\begin{align*}
 \mathscr R(U,\ell;P_U)&=\mu_{P_U}+O_\rho(L_U^{-3/2})
                                    =O_\rho(L_U^{-1}),\\
 \frac1\ell\int_U^{U+\ell}J_{1/L_X}(X)\dd X
   &=O_\rho(L_U^{-1/3}),\\
 \frac{|\{X\in[U,U+\ell]:J_{1/L_X}(X)\ge h\}|}{\ell}
   &=O_{h,\rho}(L_U^{-1})\qquad(h>0\ \text{fixed}).
\end{align*}
\end{corollary}
\begin{proof}
The first assertion uses \(\mu_P\le2P^{-1/2}\).
For the Green transfer, \(L_X=L_U+O(1)\) and \(P_U\le L_X^2\)
throughout the interval. The phase gates in
Lemma~\ref{lem:joint-phase} thus hold uniformly by
\eqref{eq:euler-phase-derivatives}. In \eqref{eq:euler-master},
the omitted primes cost \(O(P_U^{-3/2})=O(L_U^{-3})\).
The reference exterior costs \(O(L_U^{3/4}U^{-4})\) and the
complete-field logarithmic exterior \(O(U^{-3}\log^2U)\).
Since \(q_X\log b_\delta\le\delta/4\) and \(\log(1+x)\le x\),
\[
 0\le J_{1/L_X}(X)\le
     \delta/4+\frac{R_{P_U}(X,1)}{2\delta^2}+C/L_U .
\]
Taking \(\delta=L_U^{-1/3}\) and averaging proves the second
assertion. For fixed \(h>0\) take a fixed \(\delta>0\) with
\(\delta/4\le h/4\), and let \(U\) be large enough for
\(C/L_U\le h/4\). Then \(J_{1/L_X}(X)\ge h\) implies
\(R_{P_U}(X,1)\ge h\delta^2\), so Markov's inequality proves
the last assertion.
\end{proof}

The bound on total Green mass is \(O_\rho(\ell L_U^{-1/3})\),
which grows on the admissible windows. A localized defect may fit inside that
budget. A bound at every center would require further control
of \(J_\lambda(X)\) along the moving-time curve
\(\lambda=1/L_X\).
